
\begin{proposition}[The prescribed toric fibre and its conjugate over the projective line]\label{cdp:st:C.2.UN.73}
Every marked central log space classified in \ref{cdp:st:C.2.UN.71} occurs
unchanged in a smooth compact complex threefold mapping
properly to $\PP^1$, with connected complex two-torus
fibres away from two values. Its two full singular fibres
are the prescribed reduced simple-NC log space and its
complex conjugate. Both neighbourhoods have a nowhere-zero
canonical form. If its lattice index is $d=[N:\Gamma]$,
each singular support has Euler number $2d$ and the
compact total space has Euler number $4d$.

The local base coordinates at the two values, and the
maps restoring the prescribed log sections, are part
of the construction below. The assertion preserves
the specified central log spaces; it does not assert
that an arbitrary given full degeneration germ is
the germ of this compact construction.
\end{proposition}

\begin{proof}
\par\medskip\noindent\textbf{The exact initial periods and a positive symmetric part.}
Use the actual data $(N,\Gamma,\mathscr P,c)$ of \ref{cdp:st:C.2.UN.71}.
Choose $N=\Z^2$, write $\Gamma=Q\Z^2$ with
$Q\in\operatorname{Mat}_2(\Z)$ invertible over $\R$,
and choose the logarithms $C$ of its two torus
holonomies. In the vector coordinate
$(2\pi i)^{-1}Q^{-1}\log z$, the angular lattice is
\[
        N_0=Q^{-1}\Z^2\supset\Z^2,
             \qquad b=(2\pi i)^{-1}Q^{-1}C.
\tag{C.TX.1}\label{cdp:eq:C.TX.1}
\]
The two other periods of the local model are
$b_j+(\log u)/(2\pi i)e_j$. Write $b=U+iA$
with $U,A$ real matrices. Choose $R>0$ large enough
that the symmetric part $S$ of $A+RI$ is positive
definite, and put
\[
 b'=b+iRI,
 \qquad \kappa=e^{-2\pi R}>0,
 \qquad A+RI=S+aJ,
 \qquad J=\begin{pmatrix}0&1\\-1&0\end{pmatrix}.
\tag{C.TX.2}\label{cdp:eq:C.TX.2}
\]
In particular
$\exp(2\pi iQb'_j)=\kappa^{Qe_j}c_{Qe_j}$.
This changes only a recorded base scale; the
actual coordinate maps undo it at both ends.
We do not commute this imaginary shift through
the real-linear operation used next.

\par\medskip\noindent\textbf{A holomorphic family on the entire punctured line.}
Fix any integer $k\ge1$. For each constant
vector $w\in\C^2$, let
\[
 V_t(w)=w+t^kJ\overline w,
 \qquad B_t=I+t^kJ,
 \qquad 0\ne t\in\C.
\tag{C.TX.3}\label{cdp:eq:C.TX.3}
\]
This is holomorphic in $t$ for every fixed $w$.
Define the period lattice by its actual generators
\[
 \Lambda_t=V_t(N_0)
   +\sum_{j=1}^2\Z
       \left(V_t(b'_j)+\frac{\log t}{2\pi i}B_te_j\right).
\tag{C.TX.4}\label{cdp:eq:C.TX.4}
\]
Changing the logarithm by $2\pi i$ adds $B_te_j$,
which is an angular period because $\Z^2\subset N_0$.
Thus these lattices are globally single-valued
over $\C^*$, with the displayed integral change
of basis on overlaps of logarithm branches.

We prove their real rank is four everywhere.
It suffices to use $\Z^2$ in place of its finite
angular enlargement $N_0$. Subtract the real
angular combinations supplied by $U$ and by
$\operatorname{Re}(\log t/(2\pi i))$ from the
last two columns. Set
\[
 x=\operatorname{Re}(t^k),\quad
 w=\operatorname{Im}(t^k),\quad
 r=|t|^k,\quad
 y=-\frac{\log|t|}{2\pi}.
\tag{C.TX.5}\label{cdp:eq:C.TX.5}
\]
With $A'=S+aJ=A+RI$, the resulting real
period matrix is exactly
\[
 \begin{pmatrix}
 I+xJ & wJ(A'-yI)\\
 wJ   & (I-xJ)A'+y(I+xJ)
 \end{pmatrix}.
\tag{C.TX.6}\label{cdp:eq:C.TX.6}
\]
Its determinant, by the displayed block
multiplication, is
\[
 \begin{aligned}
 \Delta(t)={}&(1+r^2)^2\det S
       +(1-r^2)(1+r^2)y\operatorname{tr}S\\
       &+y^2(1-r^2)^2
       +\bigl((1+r^2)a+2xy\bigr)^2>0.
 \end{aligned}
\tag{C.TX.7}\label{cdp:eq:C.TX.7}
\]
Indeed $(1-r^2)y\ge0$, the first term is
strictly positive, and the other three
terms are nonnegative. The identity uses
$r^2=x^2+w^2$ and holds for every argument
of $t$; it also holds where $B_t$ fails
to be complex-invertible. Thus \eqref{cdp:eq:C.TX.4}
is a discrete rank-four real lattice
for every $t\in\C^*$.

The actual quotient of the holomorphic
vector bundle $\C^2\times\C^*$ by these
locally holomorphic period sections is
therefore a smooth family of connected
complex two-tori. On every compact
annulus, the real period matrices and
their inverses are uniformly bounded.
Their real fundamental parallelepipeds
give compact representatives. This
proves properness of this middle family
over the annulus, and hence over $\C^*$.

\par\medskip\noindent\textbf{The exact chart and prescribed log section at zero.}
Near zero $B_t$ is complex-invertible.
In the middle vector coordinate $v$ use
the actual exponential map
\[
                   z=\exp(2\pi iQ B_t^{-1}v).
\tag{C.TX.8}\label{cdp:eq:C.TX.8}
\]
It quotients exactly the angular lattice
$B_tN_0$, since $QB_t^{-1}B_tN_0=\Z^2$.
The other two period translations act by
\[
 z\longmapsto c^0_{Qe_j}(t)\,t^{Qe_j}z,
 \qquad
 c^0_{Qe_j}(t)=
       \exp\bigl(2\pi iQ B_t^{-1}V_t(b'_j)\bigr).
\tag{C.TX.9}\label{cdp:eq:C.TX.9}
\]
These are holomorphic nonzero vectors
near zero; extend them multiplicatively
to a homomorphism $c^0(t):\Gamma\to T_N$.
Their exact central values are
$c^0_\lambda(0)=\kappa^\lambda c_\lambda$.

Use the base coordinate $u=\kappa t$.
The same action is then
\[
 z\longmapsto d^0_\lambda(u)u^\lambda z,
 \qquad
 d^0_\lambda(u)=\kappa^{-\lambda}
                    c^0_\lambda(u/\kappa),
 \qquad d^0_\lambda(0)=c_\lambda.
\tag{C.TX.10}\label{cdp:eq:C.TX.10}
\]
This is an actual equality of holomorphic
maps, not a replacement of its limiting
constant by a numerical invariant.

Apply the cone-chart construction of
\ref{cdp:st:C.2.UN.71} to these variable holomorphic
units $d^0_\lambda(u)$. Multiplication
by them is a holomorphic torus action
on every toric chart over the same
base. Their homomorphism identity
makes the shear action exact. The
proof of freeness and properness is
unchanged: the real-linear maps
$\lambda\mapsto\log|d^0_\lambda(u)|$
are uniformly bounded as $u\to0$,
so the perturbation of the tropical
translation lattice tends uniformly
to the identity. The same finitely
many barycentric unit polydiscs give
properness at zero. Thus \eqref{cdp:eq:C.TX.8}
identifies this filling with the
middle family on a punctured disc.
Its actual central log constants in
the coordinate $u$ are exactly $c$,
by \eqref{cdp:eq:C.TX.10}, so the specified
marked log space is retained.

\par\medskip\noindent\textbf{The exact conjugate filling at infinity.}
Put $s=t^{-1}$ and change the vector
frame by $v_\infty=s^kv$. The angular
periods become
\[
 B^\infty_sN_0,
       \qquad B^\infty_s=s^kI+J.
\tag{C.TX.11}\label{cdp:eq:C.TX.11}
\]
Reverse the two nonangular period
generators. Their resulting exact
expressions are
\[
 -(s^kb'_j+J\overline{b'_j})
                +\frac{\log s}{2\pi i}B^\infty_se_j.
\tag{C.TX.12}\label{cdp:eq:C.TX.12}
\]
The integer generator reversal is
permitted because it changes neither
the subgroup nor the actual quotient.
The matrix $B^\infty_s$ is invertible
near $s=0$. The exponential coordinates
$z_\infty=\exp(2\pi iQ(B^\infty_s)^{-1}v_\infty)$
therefore give the exact deck constants
\[
 c^\infty_{Qe_j}(s)=
  \exp\bigl(-2\pi iQ(B^\infty_s)^{-1}
                  (s^kb'_j+J\overline{b'_j})\bigr).
\tag{C.TX.13}\label{cdp:eq:C.TX.13}
\]
At zero this is
\[
 c^\infty_\lambda(0)
       =\overline{c^0_\lambda(0)}
       =\kappa^\lambda\overline{c_\lambda}.
\tag{C.TX.14}\label{cdp:eq:C.TX.14}
\]
The sign in \eqref{cdp:eq:C.TX.12} is essential:
$\exp(-2\pi i\overline{Qb'})$
is the conjugate of $\exp(2\pi iQb')$.

Choose the base coordinate $v=\kappa s$
at infinity. Define
$d^\infty_\lambda(v)=\kappa^{-\lambda}
c^\infty_\lambda(v/\kappa)$.
Then $d^\infty_\lambda(0)=\overline{c_\lambda}$,
and the exact action is
$z_\infty\mapsto d^\infty_\lambda(v)v^\lambda z_\infty$.
The two base charts satisfy $uv=\kappa^2$;
they are literal holomorphic coordinates
at the two points of one projective line.
Use the same integral fan $\Sigma$ and
its variable-constant cone quotient
at this end. All its monomial chart
maps are defined over $\Z$, while its
central torus gluing constants are
the complex conjugates of the given
ones. Hence its marked central log
space is exactly the conjugate of
the prescribed one, with the stated
local base coordinate.

\par\medskip\noindent\textbf{Compact gluing and canonical neighbourhoods.}
The exponential maps \eqref{cdp:eq:C.TX.8} and
their infinity counterparts are
actual isomorphisms of the
punctured torus families, since
they first divide by precisely
the angular lattice and then
by the displayed remaining
period group. Glue the two
proper cone fillings to the
middle family along these
isomorphisms. Disjoint smaller
base discs and a compact
intermediate annulus show
directly that the glued total
space is Hausdorff and proper
over $\PP^1$. It is smooth,
and its fibres are connected.
Properness over the compact
base makes the total space
compact. There are no other
singular fibres, because the
middle family is smooth.

In both end fillings the form
$du\wedge d\log z_1\wedge d\log z_2$
(respectively with $v,z_\infty$)
is invariant under the actual
variable-constant shear action:
the additional logarithmic
derivatives are multiples of
the base differential and
vanish in this wedge. The
unimodular cone calculation
of \ref{cdp:st:C.2.UN.71} extends it to a
nowhere-zero three-form.
Thus both end neighbourhoods
are canonical-flat with the
full reduced simple-NC
fibres specified in the
statement. Finally Euler
additivity over the base,
with smooth torus Euler
number zero, gives
$e(Y)=2d+2d=4d$.
This completes the compact
occurrence construction.
\end{proof}


